% Lamina 6: arquitectura del detector (Ens-Vb-X), copia del diseno del tesista.
% Las coordenadas son los PIXELES de su imagen original (1180 x 430 px):
% x = 0.0122 cm por pixel (unos 13.75 cm de ancho), y = 0.019 cm por pixel hacia abajo.
% Para mover una caja, se cambian sus pixeles (esquina superior izquierda e inferior derecha).
% Colores tomados de la imagen (opcion B, decidida por el tesista): no son los de los detectores.
\definecolor{dAzul}{HTML}{0F4C81}
\definecolor{dNaranja}{HTML}{E67E22}
\definecolor{dMorado}{HTML}{7D3C98}
\definecolor{dVerde}{HTML}{1B9E77}
\definecolor{dRojo}{HTML}{C0392B}
\definecolor{dGris}{HTML}{4A5560}
\begin{tikzpicture}[x=0.0122cm, y=-0.019cm, font=\sffamily,
    % contenedor grande: borde de color y fondo muy claro
    grande/.style={draw=#1, fill=#1!7!white, line width=1pt, rounded corners=5pt},
    % caja interna: fondo blanco
    chica/.style={draw=#1, fill=white, line width=0.8pt, rounded corners=3pt},
    tcol/.style={font=\sffamily\fontsize{6}{7}\selectfont\bfseries, text=#1},          % titulos de columna
    tcaja/.style={font=\sffamily\fontsize{5.8}{6.8}\selectfont\bfseries, text=#1},   % titulo dentro de caja
    tnota/.style={font=\sffamily\fontsize{4.6}{5.6}\selectfont, text=dGris},                % texto gris pequeno
    flecha/.style={-{Stealth[length=4pt]}, line width=1.1pt, #1},
  ]

  % ======================= Titulos de columna =======================
  \node[tcol=dAzul]    at (119, 67) {1. ENTRADAS};
  \node[tcol=dNaranja] at (346, 67) {2. FILTRO VECINAL};
  \node[tcol=dAzul]    at (593, 67) {3. VENTANA (24\,h $\times$ 17)};
  \node[tcol=dMorado]  at (837, 67) {4. MODELO PROFUNDO};
  \node[tcol=dVerde]   at (1064, 67) {5. DECISIÓN DIARIA};

  % ======================= 1. Entradas =======================
  \draw[grande=dAzul] (30, 88) rectangle (208, 160);
  \node[tcaja=dAzul] at (119, 112) {Estación objetivo (\textit{s})};
  \node[tnota]       at (119, 138) {O\textsubscript{3} horario (ppb)};

  \draw[grande=dNaranja] (30, 183) rectangle (208, 278);
  \node[tcaja=dNaranja] at (119, 206) {4 Estaciones vecinas};
  \foreach \k [count=\n from 0] in {1,...,4} {
    \draw[chica=dNaranja] (46+\n*37, 230) rectangle (75+\n*37, 262);
    \node[font=\sffamily\fontsize{5.2}{6.2}\selectfont\bfseries, text=dNaranja] at (60.5+\n*37, 246) {v\textsubscript{\k}};
  }

  \draw[grande=dGris] (30, 302) rectangle (208, 375);
  \node[tcaja=dGris] at (119, 326) {Perfil propio};
  \node[tnota]       at (119, 352) {Últimos 14 días};

  % ======================= 2. Filtro vecinal =======================
  \draw[grande=dNaranja] (247, 138) rectangle (445, 323);
  \node[tcaja=dNaranja] at (346, 163) {Filtro ($\tau = 0.26\sigma$)};

  \draw[chica=dRojo, fill=dRojo!8!white] (257, 190) rectangle (435, 238);
  % tache dibujado (la fuente no trae el simbolo)
  \draw[dRojo, line width=1pt] (263,204) -- (270,212) (270,204) -- (263,212);
  \node[tcaja=dRojo, font=\sffamily\fontsize{5.2}{6.2}\selectfont\bfseries] at (352, 208) {Excluye sospechosas};
  \node[tnota]       at (346, 227) {(si son 3 o 4, ninguna)};

  \draw[chica=dVerde, fill=dVerde!8!white] (257, 255) rectangle (435, 303);
  % palomita dibujada
  \draw[dVerde, line width=1pt] (262,272) -- (265,277) (265,277) -- (272,266);
  \node[tcaja=dVerde, font=\sffamily\fontsize{5.2}{6.2}\selectfont\bfseries] at (352, 273) {Base vecinal limpia};
  \node[tnota]        at (346, 292) {Vecinas restantes};

  % ======================= 3. Ventana =======================
  \draw[grande=dAzul] (485, 88) rectangle (702, 375);
  \node[tcaja=dAzul] at (593, 112) {24 horas $\times$ 17 columnas};
  % \n = numero de caja, \c = color, \t = texto
  \foreach \n/\c/\t in {0/dAzul/{1. Desvío de \textit{s} vs.\ base},
                        1/dGris/{2. Hora: sin(t)},
                        2/dGris/{3. Hora: cos(t)},
                        3/dGris/{4. Máscara (¿hay dato?)},
                        4/dNaranja/{5. Desvío vs.\ vecinas},
                        5/dNaranja/{6. Igual para v\textsubscript{1}–v\textsubscript{4}}} {
    \draw[chica=\c] (502, 137+\n*38) rectangle (685, 166+\n*38);
    \node[font=\sffamily\fontsize{5.2}{6.2}\selectfont\bfseries, text=\c] at (593.5, 151.5+\n*38) {\t};
  }

  % ======================= 4. Modelo profundo =======================
  \draw[grande=dMorado] (741, 88) rectangle (933, 375);
  \node[tcaja=dMorado] at (837, 112) {Ensamble (5 LSTM)};
  \foreach \n in {0,...,4} {
    \draw[chica=dMorado] (760, 136+\n*35) rectangle (913, 163+\n*35);
    \node[font=\sffamily\fontsize{5.2}{6.2}\selectfont\bfseries, text=dMorado] at (836.5, 149.5+\n*35) {LSTM \#\the\numexpr\n+1\relax};
  }
  \draw[chica=dAzul] (760, 318) rectangle (913, 360);
  \node[font=\sffamily\fontsize{5.2}{6.2}\selectfont\bfseries, text=dAzul] at (836.5, 339) {Promedio $p_s(t)$};

  % ======================= 5. Decision diaria =======================
  \draw[grande=dVerde] (972, 88) rectangle (1157, 375);
  \node[tcaja=dVerde] at (1064, 112) {Cierre del día};

  \draw[chica=dVerde] (989, 138) rectangle (1140, 193);
  \node[tnota]        at (1064, 155) {Promedio del día};
  \node[tcaja=dVerde] at (1064, 177) {$S = \overline{p_s(t)}$};

  \draw[chica=dNaranja, fill=dNaranja!8!white] (989, 209) rectangle (1140, 266);
  \node[tcaja=dNaranja] at (1064, 226) {Umbral p99};
  \node[tcaja=texto]    at (1064, 249) {$S > 0.1508$};

  \fill[dVerde, rounded corners=4pt] (989, 283) rectangle (1140, 356);
  % Texto original del tesista. Alternativa sugerida: ALARMA DIARIA (decide al cerrar el dia).
  \node[font=\sffamily\footnotesize\bfseries, text=white, align=center] at (1064, 319.5) {ALERTA\\TEMPRANA};

  % ======================= Flechas =======================
  \draw[flecha=dAzul]    (208, 124) -- (485, 124);              % s -> ventana
  \draw[flecha=dAzul]    (208, 150) -- (247, 182);              % s -> filtro
  \draw[flecha=dNaranja] (208, 230) -- (247, 230);              % vecinas -> filtro
  \draw[flecha=dGris]    (208, 340) -- (346, 340) -- (346, 323); % perfil -> filtro (senal z)
  \draw[flecha=dVerde]   (445, 285) -- (485, 303);              % base limpia -> desvio vs vecinas
  \draw[flecha=dAzul]    (702, 231) -- (741, 231);              % ventana -> ensamble
  \draw[flecha=dVerde]   (933, 231) -- (972, 231);              % ensamble -> decision
\end{tikzpicture}